\documentclass[10pt,a4paper]{amsart}
\usepackage[margin=19mm]{geometry}
\usepackage{amsmath,amssymb,mathtools}
\usepackage[scr=rsfs,cal=euler]{mathalfa}
\usepackage{xfrac}
\usepackage[unicode,hidelinks,bookmarks=false]{hyperref}
\newcommand{\ie}{i.e.\ }
\newcommand{\cf}{cf.\ }
\newcommand{\resp}{respectively\ }
\newcommand{\Subsection}{\subsection}
\providecommand{\nobreakdash}{\nobreak\hbox{-}}
\setlength{\emergencystretch}{2em}
\multlinegap0pt
\usepackage{amssymb}
\usepackage[shortlabels]{enumitem}
\newtheorem{theorem}{Theorem}
\title{Arithmetic properties of DSOME function}
\author{Nayandeep Deka Baruah, Pankaj Gogoi}
\date{Source: arXiv:2602.20025v2 (CC BY 4.0)}
\begin{document}
\maketitle

\noindent\emph{Source and licence.} Arithmetic properties of DSOME function. Version arXiv:2602.20025v2 (CC BY 4.0), distributed by arXiv under the Creative Commons Attribution 4.0 International licence, http://creativecommons.org/licenses/by/4.0/, as recorded in the arXiv licence metadata for this exact version and shown on the abstract page https://arxiv.org/abs/2602.20025v2. Full licence notice: \texttt{licenses/arxiv-2602.20025-CC-BY-4.0.txt}.\par\smallskip

\noindent\textit{Editorial scope.} Counted unit: the article's theorem Theorem1.2 (source0.tex lines 164--169). The article's complete original proof of it is delivered with it (source0.tex lines 276--433, one \texttt{proof} environment of 158 lines, which the article places away from the statement). No mathematics is written, repaired or re-derived: the statement and the proof are the article's own lines, in the article's own notation, and no retained line is substituted or re-flowed.  Retained uncounted as the source-local context the proof needs: lines 161--162; lines 187--189; lines 195--195; lines 198--204; lines 209--214; lines 219--222; lines 228--231; lines 236--238.  Nothing was omitted from any retained span, so the package carries no omission record.  No retained line contains a citation, so the wrapper carries no bibliography and no bibliography entry.  No retained line contains a figure environment, an image-inclusion command or drawing code, so no image is delivered and the package declares no asset dependency.  Editorial formatting only, in addition: an AMS \texttt{amsart} preamble that restates the article's own declarations and macros used by the retained lines, namely the environments \texttt{theorem} on separate counters, together with the standard packages those lines need.  The retained environments print in this document as Theorem 1 in order of appearance, whereas the article prints the same items with the numbering of its own full document.  The complete native source of the article as distributed in the arXiv e-print for this exact version is bundled unchanged as \texttt{sources/195152/source0.tex}.
 \begin{align}   \sum_{n\geq0}DSOME(n)q^n=\frac{1}{8}\bigg(\frac{f_2}{f_1}-\frac{f_1^7}{f_2^3}\bigg)\label{DSOMEgf}.
 \end{align} 

\begin{align}
G(q):&=\sum_{n=0}^{\infty}\frac{q^{n^2}}{(q;q)_n}=\frac{1}{(q;\;q^5)_\infty(q^4;\;q^5)_\infty}\label{Gq}\\
\intertext{and}    H(q):&=\sum_{n=0}^{\infty}\frac{q^{n^2+n}}{(q;q)_n}=\frac{1}{(q^2;\;q^5)_\infty(q^3;\;q^5)_\infty},\label{Hq}\end{align}

\begin{align}\label{TGH}T(q):=\frac{q^{1/5}}{R(q)}=\dfrac{G(q)}{H(q)}.\end{align}

    \begin{align}
        xy^2-\frac{q^2}{xy^2}&=K,\label{Lemma3a}\\
        \frac{x^2}{y}-\frac{y}{x^2}&=\frac{4q}{K},\label{Lemma3b}\\
        \frac{y^3}{x}+q^2\frac{x}{y^3}&=K+\frac{4q^2}{K}-2q,\label{Lemma3c}\\
        x^3y+\frac{q^2}{x^3y}&=K+\frac{4q^2}{K}+2q,\label{Lemma3d}\\
        x^5-\dfrac{q^2}{x^5}&=\dfrac{f_1^6}{f_5^6}+11q,\label{x51byx5}
    \end{align}

    \begin{align}
       \label{c1} f_1&=f_{25}\left(T(q^5)-q-\frac{q^2}{T(q^5)}\right)
  \intertext{ and} 
        \label{c2}\frac{1}{f_1}&=\frac{f_{25}^5}{f_5^6}\bigg(T(q^5)^4+qT(q^5)^3+2q^2T(q^5)^2+3q^3T(q^5)+5q^4-\frac{3q^5}{T(q^5)}\\
        &\quad+\frac{2q^6}{T(q^5)^2}-\frac{q^7}{T(q^5)^3}+\frac{q^8}{T(q^5)^4}\bigg)\notag.
    \end{align}

\begin{align}
    G(q)=\frac{f_8}{f_2}\left(G\left(q^{16}\right)+qH(-q^4)\right)\label{c5}\\\intertext{and}
    H(q)=\frac{f_8}{f_2}\left(G(-q^{4})+q^3H(q^{16})\right)\label{watson}.
\end{align}

    \begin{align}
        \frac{f_5^5}{f_1^4f_{10}^3}&=\frac{f_5}{f_2^2f_{10}}+4q\frac{f_{10}^2}{f_1^3f_2},\label{Lemma4a}\\
        \frac{f_2^3f_5^2}{f_1^2f_{10}^2}&=\frac{f_5^5}{f_1f_{10}^3}+q\frac{f_{10}^2}{f_2}.\label{Lemma4b}
     \end{align}

    \begin{align}\label{lemmaG}
        G(q)G(q^4)\pm qH(q)H(q^4)\equiv \frac{1}{f_2}\pmod{2}.
    \end{align}

\begin{theorem}\label{Theorem1.2}
For all $n\geq0$, we have
\begin{align}
DSOME(25n+5r+1)&\equiv0\pmod{4}, ~where~  1\leq r\leq4.\notag
\end{align}
\end{theorem} 

\begin{proof}[Proof of Theorem \ref{Theorem1.2}] From  \eqref{DSOMEgf} and \eqref{Lemma4a}, 
 \begin{align*}
\sum_{n=0}^{\infty}DSOME(n)q^n
&=\frac{f_2}{8f_1}-\frac{1}{8}\left(\frac{f_1^3f_5^4}{f_2f_{10}^2}-4q\frac{f_1^4f_{10}^3}{f_2^2f_5}\right)\notag\\
&=\frac{f_2}{8f_1}-\frac{1}{8}\left(\frac{f_{2}f_{5}^{8}}{f_{1}f_{10}^{4}}-4qf_{5}^{3}f_{10}\right)+\frac{q}{2}\left(f_{5}^{3}f_{10}-4q\frac{f_{1}f_{10}^{6}}{f_{2}f_{5}^{2}}\right)\notag\\
&=\frac{f_2}{8f_1}-\frac{f_{2}f_{5}^{8}}{8f_{1}f_{10}^{4}}+qf_{5}^{3}f_{10}-2q^2\frac{f_{10}^{6}f_{1}}{f_{5}^{2}f_{2}},\end{align*}
which, by \eqref{c1} and \eqref{c2}, can be recast as
\begin{align}
\label{c67}&\sum_{n=0}^{\infty}DSOME(n)q^n\\
&=\frac{f_{25}^{5}f_{50}}{f_{5}^{6}}\left(\frac{1}{8}-\frac{f_{5}^{8}}{8f_{10}^{4}}\right)\left(T(q^{10})-q^2-\frac{q^4}{T(q^{10})}\right)\bigg(T^4(q^5)+qT^3(q^5)
\notag\\
&\quad+2qT^2(q^5)+3q^3T(q^5)+5q^4-3\frac{q^5}{T(q^5)}+2\frac{q^6}{T^2(q^5)}
-\frac{q^7}{T^3(q^5)}
+\frac{q^8}{T^4(q^5)}\bigg)\notag\\
&\quad+qf_{5}^{3}f_{10}-2q^2\frac{f_{25}f_{50}^{5}}{f_{5}^{2}}\left(T(q^5)-q-\frac{q^2}{T(q^5)}\right)\bigg(T^4(q^{10})+qT^3(q^{10})\notag\\
&\quad+2q^{2}T^2(q^{10})+3q^{6}T(q^{10})+5q^8-3\frac{q^{10}}{T(q^{10})}+
2\frac{q^{12}}{T^2(q^{10})}-\frac{q^{14}}{T^3(q^{10})}\notag\\
&\quad+\frac{q^{16}}{T^4(q^{10})}\bigg).\notag
\end{align}
Extracting the terms involving $q^{5n+1}$ from both sides of \eqref{c67}, dividing by $q$ and then replacing $q^5$ by $q$, we find that

\begin{align*}
&\sum_{n=0}^{\infty}DSOME(5n+1)q^n\notag\\
&=f_{1}^{3}f_{2}+\left(\frac{1}{8}-\frac{f_{1}^{8}}{8f_{2}^{4}}\right)\frac{f_{5}^{5}f_{10}}{f_{1}^{6}}\bigg(x^{3}y+\frac{q^2}{x^{3}y}-2q\left(\frac{x^2}{y}-\frac{y}{x^2}\right)
-5q\bigg)\\
&\quad-2\frac{f_{5}f_{10}^{5}}{f_{1}^{2}}\bigg(2q\left(xy^2-\frac{q^2}{xy^2}\right)-q\left(\frac{y^3}{x}+q^2\frac{x}{y^3}\right)
-5q^2\bigg),\end{align*}
which by \eqref{Lemma3a} - \eqref{Lemma3d} can be recast as
\begin{align}
\label{e1}&\sum_{n=0}^{\infty}DSOME(5n+1)q^n\\
&=f_1^3f_2+\frac{1}{8}\bigg(1-\frac{f_{1}^{8}}{8f_{2}^{4}}\bigg)\frac{f_{5}^{5}f_{10}}{f_{1}^{6}}\bigg(K-\frac{4q^2}{K}-3q\bigg)-2q\frac{f_{5}f_{10}^{5}}{f_{1}^{2}}
\bigg(K-\frac{4q^2}{K}-3q\bigg)\notag.
    \end{align}

Now, by \eqref{Lemma4a} and \eqref{Lemma4b}, we have
    \begin{align}
       \label{qq12} K-\frac{4q^2}{K}-3q&=\left(K+q\right)\left(1-\frac{4q}{K}\right)=\left(\frac{f_{2}f_{5}^{5}}{f_{1}f_{10}^{5}}+q\right)\left(1-4q\frac{f_{1}f_{10}^{5}}{f_{2}f_{5}^{5}}\right)\\
    &=\frac{f_{2}^{4}f_{5}^{2}}{f_{1}^{2}f_{10}^{4}}\times\frac{f_{1}^{4}f_{10}^{2}}{f_{2}^{2}f_{5}^{4}}=\frac{f_{1}^{2}f_{2}^{2}}{f_{5}^{2}f_{10}^{2}}.\notag
    \end{align}
  Employing \eqref{qq12} in \eqref{e1}, we have
    \begin{align}
&\sum_{n=0}^{\infty}DSOME(5n+1)q^n=f_{1}^{3}f_{2}+\frac{1}{8}\frac{f_{2}^{2}f_{5}^{3}}{f_{1}^{4}f_{10}}-\frac{1}{8}\frac{f_{1}^{4}f_{5}^{3}}{f_{2}^{2}f_{10}}-2q\frac{f_{2}^{2}f_{10}^{3}}{f_{5}}.\end{align}
Applying \eqref{Lemma4a} in the above, we find that
\begin{align}\label{c8}&\sum_{n=0}^{\infty}DSOME(5n+1)q^n\\
&=f_{1}^{3}f_{2}+\frac{1}{8}\left(\frac{f_{10}}{f_{5}}+4q\frac{f_{2}f_{10}^{4}}{f_{1}^{3}f_{5}^{2}}\right)-\frac{1}{8}\left(\frac{f_{5}^{7}}{f_{10}^{3}}-4q\frac{f_{1}f_{5}^{2}f_{10}^{2}}{f_{2}}\right)-2q\frac{f_{2}^{2}f_{10}^{3}}{f_{5}}\notag\\
&=f_{1}^{3}f_{2}+\frac{1}{8}\frac{f_{10}}{f_5}+\frac{q}{2}\left(\frac{f_{1}f_{10}^{6}}{f_{2}f_{5}^{6}}+4q\frac{f_{10}^{9}}{f_{1}^{2}f_{5}^{7}}\right)-\frac{f_{5}^{7}}{8f_{10}^{3}}+\dfrac{q}{2}\frac{f_{1}f_{5}^{2}f_{10}^{2}}{f_{2}}-2q\frac{f_{2}^{2}f_{10}^{3}}{f_{5}}\notag\\
&=\frac{f_{10}}{8f_5}-\frac{f_5^7}{8f^3_{10}}+f_1^3f_2
+\dfrac{q}{2}\left(\frac{f_{1}f_{10}^6}{f_2f_5^6}+\frac{f_{1}f_5^2f_{10}^2}{f_2}\right)-2q\frac{f_{2}^{2}f_{10}^{3}}{f_{5}}
+2q^2\frac{f_{10}^{9}}{f_{1}^{2}f_{5}^{7}}.\notag
\end{align}
Therefore, 
\begin{align*}&\sum_{n=0}^{\infty}DSOME(5n+1)q^n\\
&\equiv\frac{f_{10}}{8f_5}-\frac{f_5^7}{8f^3_{10}}+f_1^3f_2
+q\frac{f_{1}f_5^2f_{10}^2}{f_2}-2q\frac{f_{2}^{2}f_{10}^{3}}{f_{5}}
+2q^2\frac{f_{10}^{9}}{f_{1}^{2}f_{5}^{7}}\pmod{4},
\end{align*}
which by \eqref{Lemma4b} yields
\begin{align*}&\sum_{n=0}^{\infty}DSOME(5n+1)q^n\\
&=\frac{f_{10}}{8f_5}-\frac{f_5^7}{8f^3_{10}}+f_1^3f_2
+\left(\frac{f_2^3f_5^4}{f_1f_{10}^2}-\frac{f_5^7}{f_{10}^3}\right)-2q\frac{f_{2}^{2}f_{10}^{3}}{f_{5}}
+2q^2\frac{f_{10}^{9}}{f_{1}^{2}f_{5}^{7}}\notag\\
&\equiv \frac{f_{10}}{8f_5}-\frac{9}{8}\frac{f_5^7}{f_{10}^3}+2f_1f_4+2q\frac{f_4f_{10}^3}{f_5}+2q^2\frac{f_5f_{10}^5}{f_2}\pmod{4}.
\end{align*}
 Employing \eqref{c1} and \eqref{c2} in the above, we have
       \begin{align}              
    \label{c7} &\sum_{n=0}^{\infty}DSOME(5n+1)q^n\\
    &\equiv \frac{f_{10}}{8f_5}-\frac{9}{8}\frac{f_5^7}{f_{10}^3}+2f_{25}f_{100}\bigg(T(q^5)-q-\frac{q^2}{T(q^5)}\bigg)\bigg(T(q^{20})-q^4-\frac{q^8}{T(q^{20})}\bigg)\notag\\
    &\quad+2q\frac{f_{10}^3f_{100}}{f_5}
  \bigg(T(q^{20})-q^4-\frac{q^8}{T(q^{20})}\bigg)
         +2q^2\frac{f_5f_{10}^5f_{50}^5}{f_{10}^6}\bigg(T^4(q^{10})\notag\\
         &\quad+q^2T^3(q^{10})+2q^4T^2(q^{10})+3q^6T(q^{10})+5q^8-3\frac{q^{10}}{T(q^{10})}
         +2\frac{q^{12}}{T^2(q^{10})}\notag\\
         &\quad-\frac{q^{14}}{T^3(q^{10})}+\frac{q^{16}}{T^4(q^{10})}\bigg)\pmod{4}\notag.
       \end{align}
       
       To derive Theorem \ref{Theorem1.2}, we need to extract the terms involving $q^{5n+r}$ for $r=1,2,3$, and $4$ in  \eqref{c7}. We treat each of the cases for $r$ separately in the following. 

      \vskip .4cm \noindent \emph{Case $r=1$}. Extracting the terms involving $q^{5n+1}$ from both sides of \eqref{c7}, dividing by $q$, and then replacing $q^5$ by $q$, we obtain
      \begin{align*} &\sum_{n=0}^{\infty}DSOME(25n+6)q^n\\
      &\equiv 2f_5f_{20}\bigg(-T(q^4)+\frac{q}{T(q)}\bigg)+2\frac{f_2^3f_{20}T(q^4)}{f_1}+2q^3\frac{f_1f_{10}^5}{f_2T^3(q^2)}\pmod{4},\end{align*}
      which by \eqref{TGH} can be rewritten as
      \begin{align}\label{ty5}&\sum_{n=0}^{\infty}DSOME(25n+6)q^n\\
      &\equiv 2f_5f_{20}\bigg(\frac{G(q^4)}{H(q^4)}-q\frac{H(q)}{G(q)}\bigg)+2f_1f_4f_{20}\frac{G(q^4)}{H(q^4)}+2q^3\frac{f_1f_{10}^5H^3(q^2)}{f_2G^3(q^2)}\notag\\
         &\equiv 2\frac{f_5f_{20}}{G(q)H(q^4)}\bigg(G(q)G(q^4)-qH(q)H(q^4)\bigg)+2f_1f_4f_{20}\frac{G(q^4)}{H(q^4)}\notag\\&\quad+2q^3\frac{f_1f_{10}f_{40}H(q^2)H(q^4)}{f_2G(q^2)G(q^4)}\pmod{4},\notag
    \end{align}
    where we also used the fact that for any positive integer $j$,
    \begin{align}\label{HGmod2}G^2(q^j)\equiv G(q^{2j})\pmod{2}\quad\textup{and}\quad H^2(q^j)\equiv H(q^{2j})\pmod{2}.
    \end{align}

    Next, it follows from \eqref{Gq} and \eqref{Hq} that
     \begin{align} H(q)G(q)&=\frac{f_{5}}{f_1}.\label{HqGq}
\end{align}
Employing \eqref{lemmaG}, \eqref{HGmod2}, and \eqref{HqGq} in \eqref{ty5}, we find that
 \begin{align}   \label{dsome25n6}&\sum_{n=0}^{\infty}DSOME(25n+6)q^n\\
 &\equiv 2\frac{f_5f_{20}}{f_2G(q)H(q^4)}+2f_1f_8G^2(q^4)+2q^3f_1f_{40}\dfrac{H(q^8)}{G(q^4)}\notag\\ 
&\equiv 2f_1f_2H(q)G(q^4)+2f_1f_8G(q^4)\left(G(q^4)+q^3H(q^{16})\right)\pmod{4}.\notag
\end{align}

However, from \eqref{watson}, we have
       \begin{align} H(q)&\equiv \frac{f_8}{f_2}\left(G(q^4)+q^3H(q^{16})\right)\pmod{2}.\label{H(1)}
        \end{align}
Employing \eqref{H(1)} in \eqref{dsome25n6}, we find that
    \begin{align}
\sum_{n=0}^{\infty}DSOME(25n+6)q^n&\equiv 2f_1f_2H(q)G(q^4)+2f_1f_2H(q)G(q^4)\equiv0\pmod{4}.\notag 
        \end{align}
        This completes the proof of the case $r=1$.
        
      \vskip .4cm \noindent \emph{Case $r=2$}.
  Extracting the terms involving $q^{5n+2}$\ from both sides of \eqref{c7}, dividing by $q^2$, and then replacing $q^5$ by $q$, we obtain
  \begin{align*}
\sum_{n=0}^{\infty}DSOME(25n+11)q^n
    &\equiv 2f_5f_{20}\frac{T(q^4)}{T(q)}+2\frac{f_1f_{10}^5}{f_2}\left(T^4(q^2)-\frac{3q^2}{T(q^2)}\right)\pmod{4}.
    \end{align*}
    Using \eqref{TGH}, \eqref{HGmod2}, and \eqref{HqGq}, we find that
  \begin{align*}
&\sum_{n=0}^{\infty}DSOME(25n+11)q^n\\
    &\equiv 2f_5f_{20}\frac{G(q^4)H(q)}{H(q^4)(q)}+2\frac{f_1f_{10}^5}{f_2}\left(\frac{G^4(q^2)}{H^4(q^2)}+q^2\frac{H(q^2)}{G(q^2)}\right)\notag\\    &\equiv 2f_1f_4H^2(q)G^2(q^4)+2\frac{f_1f_{10}^5}{f_2}\left(\frac{G(q^8)}{H(q^8)}+q^2\frac{H(q^2)}{G(q^2)}\right)\notag\\
      &\equiv 2f_1f_4H(q^2)G(q^8)+2\frac{f_1f_{10}^5}{f_2G(q^2)H(q^8)}\left(G(q^2)G(q^8)+q^2H(q^2)H(q^8)\right)\pmod{4}.
      \end{align*}
Using \eqref{lemmaG} in the above and then employing \eqref{HqGq}, we obtain
   \begin{align}  \sum_{n=0}^{\infty}DSOME(25n+11)q^n &\equiv 2f_1f_4H(q^2)G(q^8)+2\frac{f_1f_{10}^5}{f_2f_4G(q^2)H(q^8)}\notag\\
      &\equiv 2f_1f_4H(q^2)G(q^8)+2\frac{f_1f_4f_{10}f_{40}}{f_2f_8G(q^2)H(q^8)}\notag\\ &\equiv 2f_1f_4H(q^2)G(q^8)+ 2f_1f_4H(q^2)G(q^8)\notag\\
      &\equiv  4f_1f_4H(q^2)G(q^8)\notag\\
      &\equiv 0\pmod{4},\notag
      \end{align}
which completes the proof of the case $r=2$.

      \vskip .4cm \noindent \emph{Case $r=3$}.
We extract the terms involving $q^{5n+3}$ from both sides of \eqref{c7} and then proceed as in the previous two cases. Accordingly, we deduce that
\begin{align*}
&\sum_{n=0}^{\infty}DSOME(25n+16)q^n\notag\\
&\equiv 2qf_5f_{20}\frac{T(q)}{T(q^4)}+2q\frac{f_1f_{10}^5}{f_2}\left(T(q^2)+\frac{q^2}{T^4(q^2)}\right)\notag\\
    &\equiv 2qf_5f_{20}\frac{H(q^4)G(q)}{G(q^4)H(q)}+2q\frac{f_1f_{10}f_{40}}{f_2}\left(\frac{G(q^2)}{H(q^2)}+q^2\frac{H^4(q^2)}{G^4(q^2)}\right)\notag\\
     &\equiv 2qf_5f_{20}\frac{H(q^4)G(q)}{G(q^4)H(q)}+2q\frac{f_1f_{10}f_{40}}{f_2H(q^2)G(q^8)}\left(G(q^2)G(q^8)+q^2G(q^2)G(q^8)\right)\notag\\
&\equiv 2qf_1f_4G(q^2)H(q^8)+2q\frac{f_1f_{10}f_{40}}{f_2f_4H(q^2)G(q^8)}\notag\\     
&\equiv 2qf_1f_4G(q^2)H(q^8)+2qf_1f_4G(q^2)H(q^8)\notag\\ 
&\equiv0\pmod{4}.
      \end{align*}
          
     \vskip .4cm \noindent \emph{Case $r=4$}. For this case, we extract the terms involving $q^{5n+4}$ from both sides of \eqref{c7} and then proceed as in the case for $r=1$ to deduce that     
    \begin{align}  \label{case4a} &\sum_{n=0}^{\infty}DSOME(25n+21)q^n\equiv  2f_5f_{20}\left(T(q)+\frac{q}{T(q^4)}\right)+2q\frac{f_1f_2^2f_{20}}{T(q^4)}+2\frac{f_1f_{10}^5}{f_2}T^3(q^2)\\
      &\equiv  \frac{2f_5f_{20}}{H(q)G(q^4)}\left(G(q)G(q^4)+qH(q)H(q^4)\right)+2qf_1f_4f_{20}\dfrac{H(q^4)}{G(q^4)}\notag\\
      &\quad+2\frac{f_1f_{10}f_{40}G(q^2)G(q^4)}{f_2H(q^2)H(q^4)}\notag\\
      &\equiv  2\frac{f_5f_{20}}{f_2H(q)G(q^4)}+2qf_1f_4^2H(q^8)+2f_1f_8H(q^4)G(q^{16})\notag\\
      &\equiv 2f_1f_2G(q)H(q^4)+2f_1f_8H(q^4)\left(G(q^{16})+qH(q^4)\right)\pmod{4}.\notag
       \end{align}
     
     Now, by \eqref{c5}, we have
     \begin{align}\label{Gmod2}
     G(q)\equiv \dfrac{f_8}{f_2}\left(G(q^{16})+qH(q^4)\right)\pmod{4}.
     \end{align}
     Using \eqref{Gmod2} in \eqref{case4a}, we arrive at
    \begin{align}
      \sum_{n=0}^{\infty}DSOME(25n+21)q^n
      &\equiv 4f_1f_2G(q)H(q^4)\equiv  0\pmod{4}.\notag
       \end{align}
     This completes the proof.  
  \end{proof}

\end{document}
